\section{Conservación bilateral y transducción aritmético-geométrica}
\label{lib04:bilateral}

La constante holográfica de acoplamiento aritmético-geométrico designa
el registro generado \(K\), con sus ventanas ordenadas, origen de fase,
orientación e incidencias. Su producción por el registro firmado y la
inversión de Hadamard está desarrollada en la \cref{sec:k-registro}.
La lectura racional \(\kappa\), el marco orbital \(O_K\) y los
operadores que transportan sus realizaciones son objetos diferentes.
Esta distinción permite componer sus funciones: generación del registro,
codificación aritmética, transporte de información y recuperación de las
lecturas geométricas.

Las coordenadas regionales y \(\alpha\) se reciben de la genealogía
APP--TRIT--TPK, del mismo estado enriquecido y de su estructura discreta
conjunta del continuo. Las construcciones siguientes actúan sobre esas
realizaciones posteriores. Su contenido adicional es una conservación
bilateral independiente de la dimensión, una extensión reversible que
admite memoria entrante y una política de archivo con recuperación
uniformemente controlada.

\subsection{Ley bilateral para dos transportes isométricos}
\label{lib04:ley}

Sean \(H,H'\) espacios de Hilbert y
\(\mathcal B,\mathcal C:H\to H'\) isometrías lineales, de modo que
\(\mathcal B^*\mathcal B=\mathcal C^*\mathcal C=I_H\).
La tipografía caligráfica distingue estos transportes de la base
aritmética del refinamiento. Para \(a>0\), definimos
\begin{equation}
 \begin{aligned}
 p&=a^2+a+1=a(a+1)+1,\\
 N&=(2a+1)p=(a+1)^3+a^3,\\
 Q_-&=a\mathcal B-\mathcal C,\qquad
 Q_+=(a+1)\mathcal B+\mathcal C.
 \end{aligned}
 \label{lib04:def-q}
\end{equation}

\begin{theorem}[Balance de dos realizaciones isométricas]
\label{lib04:teo-balance}
Para las aplicaciones \eqref{lib04:def-q},
\begin{equation}
 (a+1)Q_-^*Q_-+aQ_+^*Q_+=NI_H.
 \label{lib04:balance}
\end{equation}
La identidad vale en toda dimensión, incluidos espacios de Hilbert
infinitos, y admite transportes rectangulares con dominio y codominio
comunes.
\end{theorem}

\begin{proof}
Escribamos \(X=\mathcal B^*\mathcal C+\mathcal C^*\mathcal B\).
La isometría de ambos transportes da
\begin{align*}
 Q_-^*Q_-&=(a^2+1)I-aX,\\
 Q_+^*Q_+&=((a+1)^2+1)I+(a+1)X.
\end{align*}
Las ponderaciones \(a+1\) y \(a\) cancelan exactamente los términos
de \(X\). El coeficiente restante es
\[
 (a+1)(a^2+1)+a((a+1)^2+1)
 =2a^3+3a^2+3a+1=N.
\]
Las composiciones son de operadores acotados y el cálculo utiliza
únicamente sus productos y adjuntos. Por ello la prueba se aplica
también en dimensión infinita.
\end{proof}

Para \(a=8\), recibido de la partición nonádica precedente,
\begin{equation}
 9\|Q_-v\|^2+8\|Q_+v\|^2
 =(9^3+8^3)\|v\|^2=17\cdot73\,\|v\|^2.
 \label{lib04:balance-canonico}
\end{equation}
El factor \(73=8\cdot9+1\) normaliza esta identidad general. En el
módulo ternario se realiza además como determinante y factor de
semejanza; cada una de esas funciones conserva sus hipótesis propias.
El balance \eqref{lib04:balance} utiliza las isometrías y admite tanto
transportes con orden finito como transportes sin esa propiedad.

\subsection{Factorización, recuperación y compatibilidad de los canales}
\label{lib04:factorizacion}

Definamos
\begin{equation}
 E_av=\frac1{\sqrt p}
 \begin{pmatrix}\sqrt{a(a+1)}\,\mathcal Bv\\\mathcal Cv\end{pmatrix},
 \qquad
 R_a=\frac1{\sqrt{2a+1}}
 \begin{pmatrix}\sqrt a&-\sqrt{a+1}\\
                 \sqrt{a+1}&\sqrt a\end{pmatrix}.
 \label{lib04:factorizacion-er}
\end{equation}
La matriz \(R_a\) actúa en los dos canales y se entiende tensorizada
con \(I_{H'}\). Se comprueba
\(E_a^*E_a=I_H\), \(R_a^*R_a=I_{H'\oplus H'}\), y
\begin{equation}
 W_av=R_aE_av
 =\frac1{\sqrt N}
 \begin{pmatrix}\sqrt{a+1}\,Q_-v\\\sqrt a\,Q_+v\end{pmatrix},
 \qquad W_a^*W_a=I_H.
 \label{lib04:w}
\end{equation}
En el caso \(a=8\), el primer reparto tiene pesos \(72/73\) y
\(1/73\); la segunda operación es una rotación entre los canales.
Esta factorización expresa el papel de \(72+1\) dentro de la norma
total, distinguiéndolo de la restitución lineal de
\eqref{lib03:cadena-restauracion}.

\begin{proposition}[Recuperación bilateral]
\label{lib04:prop-recuperacion}
Los canales sin normalizar \(x=Q_-v\), \(y=Q_+v\) determinan
\begin{equation}
 \mathcal Bv=\frac{x+y}{2a+1},\qquad
 \mathcal Cv=\frac{ay-(a+1)x}{2a+1},\qquad
 v=\mathcal B^*\frac{x+y}{2a+1}.
 \label{lib04:inversa-canales}
\end{equation}
Desde los canales normalizados se recupera
\(v=W_a^*W_av\), con \(\|W_a^*\|=1\) cuando \(H\neq\{0\}\).
\end{proposition}

\begin{proof}
Sumar las dos definiciones de \(Q_\pm\) elimina \(\mathcal Cv\).
La combinación \(aQ_+-(a+1)Q_-\) elimina \(\mathcal Bv\).
La primera isometría tiene inversa izquierda \(\mathcal B^*\), lo
que demuestra las fórmulas. La última afirmación se sigue de
\eqref{lib04:w}: el adjunto de una isometría tiene norma uno.
\end{proof}

La imagen \(\operatorname{im}W_a\) es cerrada y su proyector
ortogonal es \(W_aW_a^*\). Un par normalizado es compatible con una
entrada única exactamente cuando pertenece a esa imagen; su componente
ortogonal cuantifica el defecto de compatibilidad.

Si \(\mathcal B=I\) y \(\mathcal C=U\) es unitario, los dos
\(Q_\pm\) son polinomios del mismo operador. Para los canales sin
normalizar, el criterio se simplifica a
\begin{equation}
 Q_+x=Q_-y.
 \label{lib04:compatibilidad-unitaria}
\end{equation}
En efecto, la igualdad equivale a
\(U(x+y)=ay-(a+1)x\). Definiendo
\(v=(x+y)/(2a+1)\), se obtiene por sustitución
\(Q_-v=x\), \(Q_+v=y\). La necesidad resulta de la conmutación
de los dos polinomios. Esta forma simplificada pertenece al dominio
unitario indicado; el criterio de imagen cubre los transportes
generales entre espacios distintos.

\subsection{Transformación de lectores y simetría relativa}
\label{lib04:lectores}

\begin{theorem}[Lector común en ambos canales]
\label{lib04:teo-lector}
Para todo operador acotado autoadjunto \(F\) sobre \(H'\),
\begin{equation}
 W_a^*\operatorname{diag}(F,F)W_a
 =\frac{a(a+1)\mathcal B^*F\mathcal B+
                  \mathcal C^*F\mathcal C}{a(a+1)+1}.
 \label{lib04:lector-general}
\end{equation}
Si \(\mathcal B\) es unitario y \(\mathcal C=\mathcal BR\), con
\(R\) unitario, la expresión, para \(G=\mathcal B^*F\mathcal B\),
es
\begin{equation}
 \mathcal T_a(G)=\frac{a(a+1)G+R^*GR}{a(a+1)+1}.
 \label{lib04:lector-relativo}
\end{equation}
Se cumple \(\mathcal T_a(G)=G\) si y sólo si \(GR=RG\).
\end{theorem}

\begin{proof}
La expansión de \((a+1)Q_-^*FQ_-+aQ_+^*FQ_+\) cancela los
términos \(\mathcal B^*F\mathcal C+\mathcal C^*F\mathcal B\).
Los coeficientes de los términos restantes son
\(a(a+1)(2a+1)\) para \(\mathcal B^*F\mathcal B\) y \(2a+1\)
para \(\mathcal C^*F\mathcal C\). Dividir por \(N\) demuestra
\eqref{lib04:lector-general}. Sustituir \(\mathcal C=\mathcal BR\)
produce \eqref{lib04:lector-relativo}. La igualdad con \(G\) equivale
a \(R^*GR=G\), y la unitariedad de \(R\) la hace equivalente a
\(GR=RG\).
\end{proof}

El lector \(F=I\) recupera la norma. Para un lector general, la
identidad determina su transformación, y la coincidencia
\(\mathcal B^*F\mathcal B=\mathcal C^*F\mathcal C=G\) garantiza
su conservación en todos los estados. En el caso canónico,
\begin{equation}
 \mathcal T_8(G)=\frac{72G+R^*GR}{73}.
 \label{lib04:lector-72}
\end{equation}
Esta fórmula relaciona conservación y simetría mediante un conmutador
concreto. El observable que transporta exactamente un operador
\(A\) del dominio a la imagen completa es \(W_aAW_a^*\); su
identidad ambiente es \(W_aW_a^*\). El observable diagonal de
\eqref{lib04:lector-general} es otra elección, cuya diferencia queda
calculada por el mismo teorema.

\subsection{Actualización reversible con memoria entrante}
\label{lib04:reentrada}

Sean ahora \(\mathcal B=I\), \(\mathcal C=U\), con \(U\)
unitario en un mismo espacio, y escribamos
\begin{equation}
 A=\frac{\sqrt{a+1}}{\sqrt N}Q_-,\qquad
 D=\frac{\sqrt a}{\sqrt N}Q_+.
 \label{lib04:ad}
\end{equation}
Los operadores \(A,D\) son normales y conmutan entre sí y con sus
adjuntos, por ser polinomios de \(U\) y \(U^*\).

\begin{theorem}[Extensión unitaria]
\label{lib04:teo-unitaria}
La actualización
\begin{equation}
 \mathscr R_a=
 \begin{pmatrix}A&-D^*\\D&A^*\end{pmatrix}:H\oplus H\to H\oplus H
 \label{lib04:unitaria}
\end{equation}
es unitaria. Su primera columna es \(W_a\), y una entrada
\((v,m)\) conserva \(\|v\|^2+\|m\|^2\).
\end{theorem}

\begin{proof}
El balance da \(A^*A+D^*D=I\). La normalidad proporciona
\(AA^*+DD^*=I\). Los productos cruzados de
\(\mathscr R_a^*\mathscr R_a\) y
\(\mathscr R_a\mathscr R_a^*\) se anulan por las conmutaciones
anteriores; sus dos diagonales son la identidad. La primera columna
reproduce \eqref{lib04:w}, y la conservación de la norma conjunta
es la identidad unitaria aplicada a \((v,m)\).
\end{proof}

Para unos datos \(z=(z_1,z_2)\), la inversión explícita es
\begin{equation}
 v_{\mathrm{rec}}=A^*z_1+D^*z_2,\qquad
 m_{\mathrm{rec}}=-Dz_1+Az_2.
 \label{lib04:inversa-unitaria}
\end{equation}
La hipótesis de memoria entrante nula se comprueba mediante
\(m_{\mathrm{rec}}=0\). Si \(e\) es una perturbación de los datos,
la unitariedad da
\begin{equation}
 \|e\|^2=\|e_{\mathrm{rec}}\|^2+\|e_{\mathrm{mem}}\|^2.
 \label{lib04:descomposicion-ruido}
\end{equation}
La segunda componente detecta la parte incompatible con esa entrada.
Un error perteneciente a \(\operatorname{im}W_a\) tiene residuo de
memoria nulo y altera la entrada recuperada. La corrección del registro
entero utiliza, además, sus distancias discretas de separación.

\subsection{Reducción causal de la actualización bilateral}
\label{lib04:reduccion-causal}

La actualización unitaria anterior realiza el transporte del estado y de su
memoria en un mismo espacio. La observación de la primera componente admite
una descripción equivalente que conserva explícitamente la contribución del
registro. Esta descripción se obtiene del transporte ya construido, con los
mismos coeficientes y la misma orientación; la memoria eliminada de la lista
de variables reaparece como dependencia de los estados anteriores.

Fijemos los operadores \(A,D\) de \eqref{lib04:ad} y una realización
\(H\oplus H\). Para \(n\geq0\), sean
\begin{equation}
 v_{n+1}=Av_n-D^*m_n,\qquad
 m_{n+1}=Dv_n+A^*m_n.
 \label{lib04:recurrencia-memoria}
\end{equation}
La actualización conjunta es \(\mathscr R_a\) y conserva
\(\|v_n\|^2+\|m_n\|^2\).

\begin{proposition}[Eliminación exacta del registro auxiliar]
\label{lib04:prop-memoria-reducida}
Para todo \(n\geq0\), con la suma vacía interpretada como cero,
\begin{align}
 m_n&=(A^*)^n m_0+
       \sum_{k=0}^{n-1}(A^*)^{n-1-k}Dv_k,
       \label{lib04:memoria-explicita}\\
 v_{n+1}&=Av_n-D^*(A^*)^n m_0
       +\sum_{k=0}^{n-1}\mathscr K_{n-1-k}v_k,
       \label{lib04:evolucion-reducida}\\
 \mathscr K_j&=-D^*(A^*)^jD\quad(j\geq0).
       \label{lib04:nucleo-retardo}
\end{align}
La ecuación reducida, junto con \eqref{lib04:memoria-explicita}, es
equivalente a la recurrencia conjunta para los mismos datos \((v_0,m_0)\).
\end{proposition}
\begin{proof}
La primera identidad vale para \(n=0\). Si vale para \(n\), la segunda
ecuación de \eqref{lib04:recurrencia-memoria} proporciona
\[
 m_{n+1}=(A^*)^{n+1}m_0+
 \sum_{k=0}^{n-1}(A^*)^{n-k}Dv_k+Dv_n,
\]
que es la identidad en \(n+1\). Sustituirla, en profundidad \(n\), en
la primera ecuación demuestra \eqref{lib04:evolucion-reducida}.
Recíprocamente, si \(v_n\) satisface esta última y se define \(m_n\)
por \eqref{lib04:memoria-explicita}, la misma diferencia de sumas prueba
\(m_{n+1}=Dv_n+A^*m_n\); la otra ecuación se recupera por sustitución.
\end{proof}

El núcleo \(\mathscr K_j\) describe una transferencia hacia el registro,
\(j\) actualizaciones dentro de él y una reentrada en la componente leída.
El término \(-D^*(A^*)^nm_0\) conserva la memoria inicial. Su supresión
corresponde al caso particular \(m_0=0\), no a una identidad del transporte
general. En \eqref{lib04:evolucion-reducida} intervienen únicamente el
estado actual y los estados anteriores. La causalidad de esta reducción
procede del orden de la recurrencia conjunta.

La normalización del balance proporciona además una estimación explícita.
Como \(Q_-=aI-U\), se tiene
\(\|A\|^2\leq r_a=(a+1)^3/((a+1)^3+a^3)<1\), mientras que
\(D^*D\leq I\). Por tanto,
\begin{equation}
 \|\mathscr K_j\|\leq r_a^{j/2},\qquad
 \sum_{j\geq0}\|\mathscr K_j\|
 \leq\frac{1}{1-\sqrt{r_a}}.
 \label{lib04:cota-nucleo-retardo}
\end{equation}
Esta cota controla el peso de la memoria en la ecuación reducida; la norma
del estado conjunto sigue conservándose por la actualización unitaria.

La proposición es la realización temporal de la reducción de Schur de
\cref{mem:prop-schur}. En las series formales
\(V(z)=\sum_{n\geq0}v_nz^n\) y
\(M(z)=\sum_{n\geq0}m_nz^n\), las recurrencias dan
\(V=v_0+zAV-zD^*M\) y \(M=m_0+zDV+zA^*M\). Eliminando \(M\),
\begin{equation}
 V(z)=\bigl[I-zA+z^2D^*(I-zA^*)^{-1}D\bigr]^{-1}
       \bigl[v_0-zD^*(I-zA^*)^{-1}m_0\bigr].
 \label{lib04:schur-bilateral}
\end{equation}
Las inversas formales existen porque sus términos constantes son identidades;
las series de los estados y el resolvente conjunto también convergen para
\(|z|<1\). La fuente inicial y el operador efectivo se transforman juntos.
Un lector que actúe también sobre \(m_n\) conserva su contribución mediante
\eqref{lib04:memoria-explicita}, conforme a \eqref{mem:schur-lector}.

La reentrada aquí descrita actualiza un registro entrante. La política de
archivo que sigue conserva, en cambio, cada complemento en una coordenada
distinta. Ambas operaciones poseen su propia ley de conservación. El núcleo
\(\mathscr K_j\), los propagadores de Green y una fuerza dimensional son
objetos de tipos diferentes: el primero elimina una componente de estado,
los segundos resuelven ecuaciones de fuente con condiciones de soporte, y
la fuerza añade una realización dinámica y sus unidades.


\subsection{Archivo no estacionario y agotamiento uniforme del terminal}
\label{lib04:archivo}

En cada etapa sean \(H_n,H_{n+1}\) espacios de Hilbert y
\(\mathcal B_n,\mathcal C_n:H_n\to H_{n+1}\) isometrías. El
parámetro \(a>0\) permanece fijo. Sean \(A_n,D_n\) los bloques
normalizados de \eqref{lib04:w}; se continúa por el primero y se
registra el segundo:
\begin{equation}
 R_0=I,\qquad R_{n+1}=A_nR_n,\qquad m_n=D_nR_nv.
 \label{lib04:politica-archivo}
\end{equation}
Esta política posee un dominio distinto de la reintroducción de
memoria de \eqref{lib04:unitaria}; cada una conserva su actualización.

\begin{lemma}[Transferencia uniforme]
\label{lib04:lem-tasa}
En toda etapa,
\begin{equation}
 \|A_n\|^2\leq r_a:={\frac{(a+1)^3}{(a+1)^3+a^3}}<1,
 \quad D_n^*D_n\geq(1-r_a)I,
 \quad\|R_n\|^2\leq r_a^n.
 \label{lib04:tasa}
\end{equation}
\end{lemma}

\begin{proof}
La desigualdad triangular da
\(\|a\mathcal B_n-\mathcal C_n\|\leq a+1\). Su normalización
proporciona la primera cota. La diferencia
\(N-(a+1)^3=a^3>0\) prueba \(r_a<1\). El balance local implica
\(D_n^*D_n=I-A_n^*A_n\), y multiplicar las cotas de los canales
continuados produce la estimación de \(R_n\).
\end{proof}

\begin{theorem}[Memoria completa a profundidad arbitraria]
\label{lib04:teo-archivo}
La aplicación finita
\begin{equation}
 Z_Nv=(R_Nv,D_0v,D_1R_1v,\ldots,D_{N-1}R_{N-1}v)
 \label{lib04:zn}
\end{equation}
es isométrica. El archivo infinito
\begin{equation}
 Z_\infty v=(D_nR_nv)_{n\geq0}
 \label{lib04:zinfty}
\end{equation}
pertenece a la suma hilbertiana de los espacios de salida y satisface
\begin{equation}
 \|v\|^2=\sum_{n\geq0}\|m_n\|^2,\qquad
 v=\sum_{n\geq0}R_n^*D_n^*m_n,
 \qquad Z_\infty^*Z_\infty=I.
 \label{lib04:recuperacion-infinita}
\end{equation}
La serie inversora converge en norma y el adjunto tiene norma uno.
\end{theorem}

\begin{proof}
La identidad \(A_n^*A_n+D_n^*D_n=I\), conjugada por \(R_n\),
produce
\[
 R_n^*R_n-R_{n+1}^*R_{n+1}=R_n^*D_n^*D_nR_n.
\]
Sumar desde cero hasta \(N-1\) da
\begin{equation}
 I=R_N^*R_N+\sum_{n<N}R_n^*D_n^*D_nR_n.
 \label{lib04:telescopia}
\end{equation}
Esta es la identidad de Gram de \(Z_N\). Por
\eqref{lib04:tasa}, el primer término tiende a cero en norma
operatoria. La identidad límite prueba la convergencia cuadrática del
archivo y su isometría. El adjunto de esa isometría es acotado; las
truncaciones de todo archivo cuadrado-sumable convergen en la suma
hilbertiana, y su imagen por el adjunto prueba la convergencia de la
serie de \eqref{lib04:recuperacion-infinita}. Aplicarla a los datos
exactos recupera \(v\).
\end{proof}

Con sólo los primeros \(N\) complementos, el adjunto truncado tiene
residuo exacto
\begin{equation}
 v-\sum_{n<N}R_n^*D_n^*m_n=R_N^*R_Nv,
 \qquad
 \left\|v-\sum_{n<N}R_n^*D_n^*m_n\right\|
 \leq r_a^N\|v\|.
 \label{lib04:sesgo-truncado}
\end{equation}
Si el error conjunto de los datos es a lo sumo \(\varepsilon\), la
cota total es \(\varepsilon+r_a^N\|v\|\). Con \(Z_N\), incluido
su terminal, o con \(Z_\infty\), el adjunto recupera con error a lo
sumo \(\varepsilon\).

La política anterior agota el terminal por la cota uniforme probada.
Las evoluciones precedentes que poseen un terminal positivo
\(A_\infty\neq0\) conservan en su archivo la componente
\(A_\infty^{1/2}v\). La profundidad arbitraria y la eliminación del
terminal son, por tanto, conclusiones diferentes y mantienen sus
hipótesis respectivas.

\begin{proposition}[Inversa de las memorias finitas sin terminal]
\label{lib04:prop-inversa-finita}
Para \(N\geq1\), sea \(M_Nv=(D_nR_nv)_{n<N}\). Entonces
\begin{equation}
 M_N^*M_N=I-R_N^*R_N\geq(1-r_a^N)I,
 \label{lib04:gram-finito}
\end{equation}
y la inversa de mínimos cuadrados
\begin{equation}
 L_N=(I-R_N^*R_N)^{-1}M_N^*,\qquad
 L_NM_N=I,\qquad
 \|L_N\|\leq(1-r_a^N)^{-1/2}
 \label{lib04:inversa-finita}
\end{equation}
recupera exactamente la entrada con datos exactos.
\end{proposition}

\begin{proof}
La identidad de Gram es \eqref{lib04:telescopia}; la cota procede de
\eqref{lib04:tasa}. El operador positivo es invertible por su cota
inferior estricta. La composición demuestra \(L_NM_N=I\), y
\(L_NL_N^*=(M_N^*M_N)^{-1}\) prueba la norma declarada.
\end{proof}

Esta inversa finita elimina el sesgo de
\eqref{lib04:sesgo-truncado}; su amplificación de ruido queda
especificada por \eqref{lib04:inversa-finita}. El adjunto truncado,
la inversa finita y el adjunto del archivo completo son así tres
operadores con comportamientos diferenciados.

Para \(a=8\), la cota general es \(r_a=729/1241\). En el caso
\(\mathcal B=I\), \(\mathcal C=S^4\), con \(S\) el ciclo de
doce posiciones, el proyector fijo
\(P_0=(I+S^4+S^8)/3\) da
\begin{equation}
 A^*A=\frac{a+1}{(2a+1)p}
       \bigl((a-1)^2P_0+p(I-P_0)\bigr).
 \label{lib04:gram-tasa-ternaria}
\end{equation}
Como \(p-(a-1)^2=3a>0\), su norma es
\((a+1)/(2a+1)\). En \(a=8\) resulta \(9/17\); el valor propio
del sector fijo es \(441/1241\). Estas tasas pertenecen a los
transportes y a la política declarados.

Si se refinan ambos canales en un árbol, cada nivel completo es también
isométrico por composición. Sumar las normas de todas las capas
completas contaría repetidamente la misma norma inicial. La conservación
se formula mediante las aplicaciones de refinamiento entre niveles o
mediante el archivo de complementos de \eqref{lib04:zinfty}.

\subsection{Formas geométricas y conservación de su transporte}
\label{lib04:formas}

Sean \(b,b'\) formas no degeneradas, simétricas o alternantes, sobre
los espacios correspondientes, y supóngase
\(\mathcal B^*b'\mathcal B=\mathcal C^*b'\mathcal C=b\).
El mismo desarrollo bilineal demuestra
\begin{equation}
 (a+1)Q_-^*b'Q_-+aQ_+^*b'Q_+=Nb.
 \label{lib04:balance-formas}
\end{equation}
En efecto, los términos mixtos
\(\mathcal B^*b'\mathcal C+\mathcal C^*b'\mathcal B\) se cancelan
con las mismas ponderaciones; los términos puros suman \(Nb\).
Así se conserva la forma transportada sin exigir que sea positiva.
Las cotas contractivas y de ruido de la subsección anterior utilizan,
en cambio, normas de Hilbert positivas. Una forma lorentziana o
alternante conserva \eqref{lib04:balance-formas}, con sus propias
interpretaciones geométricas, y no sustituye la positividad en esas
desigualdades.

\subsection{Correlación entre canales y ensayo inverso de una lectura}
\label{lib04:correlacion}

Para \(v\neq0\), definamos
\begin{equation}
 \chi(v)=\frac{\operatorname{Re}\langle\mathcal Bv,\mathcal Cv\rangle}
                      {\|v\|^2},\qquad -1\leq\chi(v)\leq1.
 \label{lib04:chi}
\end{equation}
La cota es Cauchy--Schwarz aplicada a las dos imágenes isométricas.
Las lecturas normalizadas son
\begin{equation}
 s_-=a^2+1-2a\chi,\qquad
 s_+=(a+1)^2+1+2(a+1)\chi.
 \label{lib04:s-pm}
\end{equation}
Para \(a=8\),
\begin{equation}
 s_-=65-16\chi,\qquad s_+=82+18\chi,\qquad
 9s_-+8s_+=1241.
 \label{lib04:s-canonico}
\end{equation}
Las dos normas individuales contienen una correlación real común. El
balance es una relación exacta entre ellas, no dos determinaciones
aritméticas independientes de una constante.

Usar la salida previa \(\alpha_A\) en el ensayo inverso
\(s_-=10^4\alpha_A\) determina la lectura complementaria
\begin{equation}
 s_+=73+\frac98(73-10^4\alpha_A)
       \simeq73{,}029783595557.
 \label{lib04:ensayo-alpha}
\end{equation}
La igualdad se obtiene sustituyendo en
\eqref{lib04:s-canonico}. El promedio ponderado
\((9s_-+8s_+)/17\) permanece exactamente igual a \(73\).
La comparación conserva la genealogía previa de \(\alpha_A\): es
una condición del ensayo inverso, no otro productor independiente de
esa salida. Un cambio unitario puede alterar \(\chi\) y las dos
lecturas manteniendo el balance, aunque cambie una relación particular
de orden ternario.

\subsection{Marco cíclico y transporte de las respuestas}
\label{lib04:marco}

Sobre el portador del registro generado, sean
\(H=\mathbb R^{12}\), \((Sx)_i=x_{i+1}\), con índices cíclicos,
y
\begin{equation}
 O_K=(K,SK,\ldots,S^{11}K),\qquad
 589609I\leq O_K^*O_K\leq39225169I.
 \label{lib04:gram-k}
\end{equation}
Las cotas son las del marco cíclico del registro canónico,
demostradas en la \cref{lib01:marco-ciclico}.
Sea \(Z\) una de las isometrías completas
\(Z_N,Z_\infty\), o una composición con los archivos precedentes
que conserve sus terminales y complementos.

\begin{theorem}[Conservación del condicionamiento del marco]
\label{lib04:teo-marco}
Si \(x=O_K\beta\), el dato \(z=Zx\) determina
\begin{equation}
 \beta=O_K^{-1}Z^*z,\qquad
 \|\delta\beta\|\leq\frac{\|\delta z\|}{\sqrt{589609}}.
 \label{lib04:inversa-marco}
\end{equation}
Las doce columnas de \(ZO_K\) forman una base de su imagen, incluso
cuando el archivo tenga infinitas coordenadas.
\end{theorem}

\begin{proof}
La igualdad \(Z^*Z=I\) da
\((ZO_K)^*(ZO_K)=O_K^*O_K\). Las cotas de Gram se conservan y
\(O_K\) es invertible por la cota inferior estricta. Componer sus
inversas demuestra la recuperación, y la norma
\(\|O_K^{-1}\|\leq589609^{-1/2}\), junto con
\(\|Z^*\|=1\), da la estabilidad. El número de columnas y su
independencia determinan la dimensión de la imagen, sin identificarla
con todo el espacio de almacenamiento.
\end{proof}

Para un operador \(H_0\) que conmute con \(S\), la respuesta
vectorial completa \(y=H_0K\) determina
\begin{equation}
 H_0=O_yO_K^{-1}.
 \label{lib04:identificacion-filtro}
\end{equation}
En efecto, la conmutación implica
\(H_0O_K=(y,Sy,\ldots,S^{11}y)=O_y\). Si se recibe únicamente
\(z=Zy+e\), se define \(\widehat y=Z^*z\) y
\(\widehat H_0=O_{\widehat y}O_K^{-1}\). Entonces
\begin{equation}
 \|\widehat H_0-H_0\|
 \leq\sqrt{\frac{12}{589609}}\,\|e\|.
 \label{lib04:error-filtro}
\end{equation}
La matriz orbital de \(\delta y\) tiene doce columnas de norma
\(\|\delta y\|\), por lo que su norma operatoria no supera
\(\sqrt{12}\|\delta y\|\). Esto y
\eqref{lib04:gram-k} prueban \eqref{lib04:error-filtro}.
Para un operador general el protocolo requiere las doce respuestas
\(H_0S^jK\); una respuesta en el caso cíclico significa sus doce
componentes, no una medición escalar.

En la imagen transportada, el desplazamiento es
\(S_Z=ZSZ^*\), y satisface \(S_Z^jZ=ZS^j\). Esta igualdad
determina el avance de fase en el archivo; una traslación de casillas
del espacio ambiente sólo representa ese avance cuando coincide con
el operador transportado.

\subsection{Lectura posicional y precisión escalar}
\label{lib04:kappa}

Fijemos \(B_0=1000\), \(d=B_0^{12}-1\), y la fila
\begin{equation}
 \ell=\frac{(B_0^{11},B_0^{10},\ldots,1)}{d}.
 \label{lib04:fila-posicional}
\end{equation}
Su extensión lineal actúa sobre \(H\), y en el dominio de palabras
\(\kappa(K)=\ell K\). El representante del lector transportado es
\(Z\ell^*\), por lo que
\begin{equation}
 \kappa(K)=\langle Z\ell^*,ZK\rangle,\qquad
 |\delta\kappa|\leq\|\ell\|\,\|e\|.
 \label{lib04:lectura-kappa}
\end{equation}
La suma geométrica proporciona
\begin{equation}
 \|\ell\|^2
 =\frac{B_0^{12}+1}{(B_0^2-1)(B_0^{12}-1)}.
 \label{lib04:norma-fila}
\end{equation}
En efecto, el numerador de la suma de cuadrados es
\((B_0^{24}-1)/(B_0^2-1)\), que al dividirse por
\((B_0^{12}-1)^2\) da \eqref{lib04:norma-fila}.

Con sólo \(N\) complementos y el adjunto truncado, aparece además
un error a lo sumo \(\|\ell\|r_a^N\|K\|\). La inversa finita
\eqref{lib04:inversa-finita} elimina ese sesgo y conserva su cota
explícita de amplificación de ruido.

En el dominio \(\{0,\ldots,999\}^{12}\), la lectura exacta
recupera \(N_K=d\kappa\) y sus doce bloques mediante divisiones
euclídeas sucesivas. Si
\begin{equation}
 |\widetilde\kappa-\kappa|<\frac1{2d},
 \label{lib04:radio-escalar}
\end{equation}
el redondeo de \(d\widetilde\kappa\) al entero más próximo
recupera \(N_K\): su distancia al entero verdadero es menor que
media unidad. La base, longitud, origen y orientación acompañan a
esa representación. El umbral escalar de
\eqref{lib04:radio-escalar} se distingue de las cotas vectoriales
de los registros de memoria.

\subsection{Incidencia, carga y restitución exacta del registro}
\label{lib04:incidencia}

Sea \(\mathbf1=(1,\ldots,1)^T\) y
\(\Pi=I-\mathbf1\mathbf1^*/12\). La isometría incidencial de
la \cref{nuc05:isometria} es
\begin{equation}
 Fv=(\mu(v),\mathcal I\Pi v/6),\qquad
 \mu(v)=\frac{\mathbf1^*v}{12},\qquad
 \|(\mu,y)\|_Y^2=12\mu^2+\|y\|^2.
 \label{lib04:f-incidencia}
\end{equation}
Su codominio es su imagen declarada. El Gram
\(\mathcal I^*\mathcal I=36I+30\mathbf1\mathbf1^*\) da
\(\|\mathcal I\Pi v\|^2=36\|\Pi v\|^2\), y
\(12\mu(v)^2+\|\Pi v\|^2=\|v\|^2\), de donde \(F^*F=I\).
Así, a partir de \(z=Zv+e\),
\begin{equation}
 \widehat{Fv}=FZ^*z,\qquad \|FZ^*e\|_Y\leq\|e\|.
 \label{lib04:error-incidencia}
\end{equation}
La incidencia sin normalizar \(\mathcal I\Pi v\) tiene error a
lo sumo \(6\|e\|\); el registro firmado \(H_{12}v\), cuya
normalización satisface \(H_{12}^*H_{12}=4I\), tiene error a lo sumo
\(2\|e\|\).

La dirección \(u=P_3\Pi K\) y el proyector
\(P_{10}=\Pi-uu^*/\|u\|^2\), con \(u\neq0\), se recuperan
con sus lectores y coordenatización previos. La selección absoluta
\(\{j:K_j\geq729\}\) tiene una discontinuidad en \(K_4=729\).
La restitución entera se efectúa antes de aplicar ese umbral.

\begin{proposition}[Decodificación después del archivo completo]
\label{lib04:prop-redondeo}
Sea \(K\) entero y \(z=ZK+e\), con \(Z^*Z=I\). Si
\(\|e\|<1/2\), redondear cada coordenada de \(Z^*z\) recupera
\(K\) exactamente. Por composición se restituyen \(\kappa\), las
firmas y los lectores geométricos definidos sobre ese registro.
\end{proposition}

\begin{proof}
La contractividad de \(Z^*\) da
\(\|Z^*z-K\|\leq\|e\|<1/2\). Cada error coordenado queda
acotado por la norma conjunta, por lo que la única coordenada entera
a distancia menor que media unidad es la original. Una vez
recuperado \(K\), evaluar cualquiera de sus lectores exactos
restituye su valor.
\end{proof}

En particular, se recupera el soporte \(\{4,6,9,10\}\), incluida
su componente de frontera. Las otras marcas que intervienen en una
bandera excepcional se conservan por sus canales propios: la
restitución de \(K\) corrige exactamente la información que tiene
en él sus argumentos suficientes.

\subsection{Conservación del condicionamiento de dos memorias}
\label{lib04:dos-memorias}

Mantengamos los lectores antecedentes
\begin{equation}
 T_j=\frac{8I+S^j}{9},\qquad
 D_j=\frac{\sqrt8}{9}(S^j-I),\qquad
 Mv=(D_3v,D_4T_3v).
 \label{lib04:m-antecedente}
\end{equation}
Los \(D_j\) de esta ecuación son los complementos precedentes;
el segundo bloque bilateral \(D\) de \eqref{lib04:ad} conserva
su propia definición. Se tiene \(\ker M=\mathbb R\mathbf1\)
y una inversa \(L_M\) de norma \(9/4\) sobre el sector centrado,
como se demuestra en la \cref{nuclear:11}.

Sea \(E=W_a\oplus W_a\) en el espacio de las dos memorias.
Su isometría implica
\begin{equation}
 (EM)^*(EM)=M^*M,\qquad
 \ker(EM)=\mathbb R\mathbf1.
 \label{lib04:gram-m-preservado}
\end{equation}
Para \(z=EMv+e\) y carga exacta
\(q=\mathbf1^*v\), la reconstrucción es
\begin{equation}
 \widehat v=\frac q{12}\mathbf1+L_ME^*z,\qquad
 \|\widehat v-v\|\leq\frac94\|e\|.
 \label{lib04:inversa-m}
\end{equation}
La prueba consiste en aplicar \(E^*E=I\), usar
\(L_MM=\Pi\) y \(\|E^*\|=1\). Si la carga tiene error
\(\delta q\), los sectores uniforme y centrado son ortogonales y
\begin{equation}
 \|\widehat v-v\|^2
 \leq\frac{|\delta q|^2}{12}+\frac{81}{16}\|e\|^2.
 \label{lib04:error-carga}
\end{equation}

Las distancias de todas las palabras de la imagen entera se conservan
exactamente por \(E\). Por tanto, los radios de restitución de
la \cref{lib02:radios-memoria} permanecen
\begin{equation}
 r_0=\frac{2\sqrt{146}}{81}\quad\hbox{sin carga},\qquad
 r_q=\frac{4\sqrt{73}}{81}\quad\hbox{con carga exacta}.
 \label{lib04:radios}
\end{equation}
Dentro del primer radio se recupera \(\Pi K\), y con ello la
dirección y su proyector cuando están definidos. Dentro del segundo,
conservando \(q\), se recuperan \(K\) y el selector absoluto de
umbral. Componer después otra isometría completa, finita o infinita,
mantiene las mismas garantías: se aplica su adjunto antes del
decodificador. El radio siempre se expresa en la norma conjunta
de los datos con la normalización fijada.

Esta composición redistribuye la información de \(M\) sin deteriorar
su condicionamiento y sin crear el canal uniforme que \(M\) elimina.
La conservación de Gram distingue esta propiedad de una mejora
numérica obtenida cambiando el lector o la escala del ruido.

\subsection{Norma orbital y normalización polar del marco}
\label{lib04:normalizacion}

La inversión precedente cumple
\(K=H_{12}U_{\mathrm{sgn}}/4\),
\(H_{12}^*H_{12}=4I\). Por tanto
\begin{equation}
 \|K\|^2=\frac14\|U_{\mathrm{sgn}}\|^2=4251257,
 \qquad \|U_{\mathrm{sgn}}\|^2=17005028.
 \label{lib04:norma-k}
\end{equation}
Cada traslación \(S^j\) es ortogonal y conserva
\(\nu=\sqrt{4251257}\). De ahí
\begin{equation}
 \operatorname{tr}(O_K^*O_K)=12\nu^2=51015084.
 \label{lib04:traza-marco}
\end{equation}
La división por \(\nu\) normaliza cada columna orbital a norma
uno y conserva su orden. Es uniforme en la órbita y bajo sus
transportes isométricos. Los registros distintos mantienen sus normas
propias, y las columnas normalizadas continúan teniendo un Gram
no escalar: sus valores propios son los anteriores divididos por
\(4251257\).

La normalización operatoria utiliza el Gram completo:
\begin{equation}
 G_K=O_K^*O_K,\qquad
 F_K=O_KG_K^{-1/2},\qquad
 F_K^*F_K=I.
 \label{lib04:polar}
\end{equation}
La raíz inversa existe por \eqref{lib04:gram-k}. La igualdad se
comprueba multiplicando por \(G_K^{-1/2}\) a ambos lados del Gram.
Así \(ZF_K\) es isométrico a cualquier profundidad.
La recuperación íntegra del marco es
\begin{equation}
 O_K=F_KG_K^{1/2}.
 \label{lib04:polar-inversa}
\end{equation}
El par \((F_K,G_K)\) conserva la información completa.
El factor unitario conserva un marco ortonormal orientado en las
coordenatizaciones fijadas, mientras que las amplitudes residen también en
\(G_K\). Multiplicar \(K\) por un escalar positivo mantiene
\(F_K\) y cambia \(G_K\); \(K\) y \(-K\) comparten
Gram y poseen factores polares opuestos, con distinta lectura lineal
\(\kappa\). Cada ejemplo muestra por qué un factor aislado no
sustituye al par.

La traslación del origen conserva \(\nu\), pero en la coordenatización
posicional fija
\begin{equation}
 \kappa(SK)=1000\kappa(K)-K_1.
 \label{lib04:cambio-origen}
\end{equation}
Se obtiene desplazando el numerador posicional y usando
\(d=1000^{12}-1\). Para conservar la misma lectura bajo un cambio
de coordenatización se transporta también la fila \(\ell\) por la inversa
de ese cambio. Una coordenatización unitaria general puede dejar la red de
bloques enteros; la decodificación por divisiones euclídeas se realiza
en la coordenatización de bloques recuperada.

\subsection{Ejemplo exacto de transducción y corrección de errores}
\label{lib04:ejemplo}

Tomemos \(a=8\), \(U=S^4\) y el registro canónico
\begin{equation}
 \begin{split}
 K=(&234,543,140,729,659,824,\\
    &621,58,914,794,146,601).
 \end{split}
 \label{lib04:k-ejemplo}
\end{equation}
Indicando por \((v)_{1:6}\) y \((v)_{7:12}\) sus dos bloques de
seis coordenadas, las lecturas sin normalizar son
\begin{align}
 (Q_-K)_{1:6}&=(1213,3520,499,5774,4358,5798),\nonumber\\
 (Q_-K)_{7:12}&=(4822,-137,7078,5809,1028,4079),\nonumber\\
 (Q_+K)_{1:6}&=(2765,5711,1881,6619,6845,8210),\nonumber\\
 (Q_+K)_{7:12}&=(5735,1123,8460,7689,1454,6138).
 \label{lib04:canales-ejemplo}
\end{align}
Su suma, dividida por \(17\), recupera todos los bloques de \(K\).
La identidad matricial
\(9Q_-^*Q_-+8Q_+^*Q_+=1241I\) antecede a su evaluación en el
vector concreto.

Para introducir un error controlado, se racionalizan las dos memorias
como \(z_0=MK/\sqrt8\), y se añade \(1/10\) a la primera
coordenada. Cada bloque de doce componentes se transporta después
por los dos canales bilaterales. La inversa por suma recupera los
datos de memoria perturbados, y el decodificador de carga recibe
\(q=6263\) junto con esos datos, sin recibir el registro buscado.
La norma del error en las memorias originales es
\(\sqrt8/10<r_q\). La normalización isométrica conserva esa norma;
los canales sin normalizar sólo facilitan el cálculo racional exacto.

En esta instancia, los candidatos de la reconstrucción por suelo y
techo que satisfacen la carga son \(924\). El mínimo de distancia
cuadrática racionalizada es \(1/100\), y el decodificador recupera
\eqref{lib04:k-ejemplo}. La unicidad se sigue del radio estricto
\eqref{lib04:radios}, no del mero tamaño de la enumeración.
La lectura aritmética recuperada es
\begin{equation}
 \kappa=
 \frac{234543140729659824621058914794146601}
      {999999999999999999999999999999999999},
 \label{lib04:kappa-ejemplo}
\end{equation}
y la selección incidencial vuelve a ser \(\{4,6,9,10\}\), con
el mismo umbral \(729\). El ejemplo introduce el ruido antes del
transporte. Las pruebas de contractividad del adjunto cubren también
errores posteriores arbitrarios cuya norma conjunta pertenezca al
radio declarado.

La capacidad operativa reunida es transportar y recuperar un registro
generado, mantener el condicionamiento de sus respuestas y restituir
sus lecturas aritméticas e incidenciales antes de aplicar selecciones
discontinuas. Sus premisas efectivas son la producción previa de
\(K\), el marco cíclico invertible, el balance bilateral, la
separación discreta de las memorias y el teorema de archivo. El
registro de doce componentes conserva el alcance de sus lectores;
la historia enriquecida completa mantiene sus otros registros y
dependencias.
